% --- Archon physics formalization source begin ---
% archon:physics
% archon:covers PhyXMiniProblems/problem_phyx_mini_0576.lean
% archon:source-report reports/phyx_mini/problem_phyx_mini_0576.source.json

\chapter{Physics problem phyx\_mini\_0576}
\label{ch:PhyXMiniProblems_problem_phyx_mini_0576}

\paragraph{Problem source.}
\#\# Physical scenario

Figure shows the decay of parents in a radioactive sample. The axes are scaled by \$N\_s = 2.00 \textbackslash{}times 10\textasciicircum{}6\$ and \$t\_s = 10.0 \textbackslash{}, \textbackslash{}text\{s\}\$.

\#\# Question

What is the activity of the sample at \$t = 27.0 \textbackslash{}, \textbackslash{}text\{s\}\$?

\#\# Answer choices

- A: A: 30Bq

- B: B: 120Bq

- C: C: 60Bq

- D: D: 6Bq

\#\# Figure caption (auxiliary; use the image as primary evidence)

The image displays a graph on a grid, featuring a curve that exhibits a decreasing trend. Here are the detailed components:

- **Axes**: The graph is plotted with two axes. The horizontal axis (x-axis) is labeled as \textbackslash{}( t \textbackslash{}) and shows the starting point as 0, ending at \textbackslash{}( t\_s \textbackslash{}). The vertical axis (y-axis) is labeled as \textbackslash{}( N \textbackslash{}) and starts at \textbackslash{}( N\_s \textbackslash{}).

- **Curve**: The curve appears to be exponential decay, starting from the top left and descending towards the bottom right. It begins at the point \textbackslash{}( (0, N\_s) \textbackslash{}) and approaches the x-axis as it moves rightward, near \textbackslash{}( t\_s \textbackslash{}).

- **Grid**: The background consists of a grid, which helps to clearly show the slope and progression of the curve.

- **Color**: The curve is a shade of magenta, providing a vivid contrast against the grid background.

This graph likely represents a concept related to exponential decay over time.

\#\# Dataset metadata

Nuclear Physics; Physical Model Grounding Reasoning, Predictive Reasoning

\paragraph{Recorded answer/context.}
C: C: 60Bq

\paragraph{Figure/image path.}
/root/proposal\_for\_physic/science-mango/phyx\_mini\_run/phyx\_data/test\_image/576.png

\paragraph{Formalization target.}
create a compiling Lean file with sorry bodies at `PhyXMiniProblems/problem\_phyx\_mini\_0576.lean`. The Lean declarations must preserve the physical quantities, dimensions or dimensional roles, figure labels, governing-law hypotheses, and final relation expressed by this problem.
Use Mathlib/Physlib names found through LeanExplore where available. If a physics API is missing, introduce faithful local abstractions rather than scalar placeholder aliases.

\begin{theorem}[Physics formalization target]
\label{thm:physics:phyx_mini_0576:target}
\lean{PhyXMiniProblems.ProblemPhyXMini0576.activity_at_twenty_seven_seconds_rounds_to_sixty_becquerels}
\uses{def:physics:phyx-mini-0576:phyxminiproblems-problemphyxmini0576-activityinbecquerels, def:physics:phyx-mini-0576:phyxminiproblems-problemphyxmini0576-radioactivedecaysetup, def:physics:phyx-mini-0576:phyxminiproblems-problemphyxmini0576-matchesproblemreadouts, def:physics:phyx-mini-0576:phyxminiproblems-problemphyxmini0576-matchessuppliedparentdecayfigure, def:physics:phyx-mini-0576:phyxminiproblems-problemphyxmini0576-hasphysicalradioactiveparameters, def:physics:phyx-mini-0576:phyxminiproblems-problemphyxmini0576-satisfiesradioactivedecaylaws, def:physics:phyx-mini-0576:phyxminiproblems-problemphyxmini0576-wholebecquerelroundingtolerance, def:physics:phyx-mini-0576:phyxminiproblems-problemphyxmini0576-iscorrectroundedanswerchoice}
The assigned autoformalize agent should translate this physics problem into Lean declarations in the covered file, with theorem and lemma proof bodies written as `by sorry`.
\end{theorem}
\begin{proof}
This is an autoformalization task, not a proof task. Produce statements that can later be proved by the physics prover without weakening the physical model.
\end{proof}

% --- Archon declaration topology begin ---
\section{Declaration topology}
\begin{definition}[Time Quantity]
  \label{def:physics:phyx-mini-0576:phyxminiproblems-problemphyxmini0576-timequantity}
  \lean{PhyXMiniProblems.ProblemPhyXMini0576.TimeQuantity}
  A nonnegative, unit-independent physical elapsed time
\end{definition}
\begin{proof}
  This is the declared model definition of Time Quantity.
\end{proof}
\begin{definition}[Decay Constant Quantity]
  \label{def:physics:phyx-mini-0576:phyxminiproblems-problemphyxmini0576-decayconstantquantity}
  \lean{PhyXMiniProblems.ProblemPhyXMini0576.DecayConstantQuantity}
  A nonnegative radioactive decay constant, carrying inverse-time dimension
\end{definition}
\begin{proof}
  This is the declared model definition of Decay Constant Quantity.
\end{proof}
\begin{definition}[Activity Quantity]
  \label{def:physics:phyx-mini-0576:phyxminiproblems-problemphyxmini0576-activityquantity}
  \lean{PhyXMiniProblems.ProblemPhyXMini0576.ActivityQuantity}
  A nonnegative radioactive activity, carrying inverse-time dimension
\end{definition}
\begin{proof}
  This is the declared model definition of Activity Quantity.
\end{proof}
\begin{definition}[time Readout]
  \label{def:physics:phyx-mini-0576:phyxminiproblems-problemphyxmini0576-timereadout}
  \lean{PhyXMiniProblems.ProblemPhyXMini0576.timeReadout}
  \uses{def:physics:phyx-mini-0576:phyxminiproblems-problemphyxmini0576-timequantity}
  Read a physical elapsed time in a selected time unit
\end{definition}
\begin{proof}
  This is the declared model definition of time Readout.
\end{proof}
\begin{definition}[decay Constant Readout]
  \label{def:physics:phyx-mini-0576:phyxminiproblems-problemphyxmini0576-decayconstantreadout}
  \lean{PhyXMiniProblems.ProblemPhyXMini0576.decayConstantReadout}
  \uses{def:physics:phyx-mini-0576:phyxminiproblems-problemphyxmini0576-decayconstantquantity}
  Read a decay constant in inverse units of a selected time unit
\end{definition}
\begin{proof}
  This is the declared model definition of decay Constant Readout.
\end{proof}
\begin{definition}[activity Readout]
  \label{def:physics:phyx-mini-0576:phyxminiproblems-problemphyxmini0576-activityreadout}
  \lean{PhyXMiniProblems.ProblemPhyXMini0576.activityReadout}
  \uses{def:physics:phyx-mini-0576:phyxminiproblems-problemphyxmini0576-activityquantity}
  Read an activity as decays per selected time unit
\end{definition}
\begin{proof}
  This is the declared model definition of activity Readout.
\end{proof}
\begin{definition}[time In Seconds]
  \label{def:physics:phyx-mini-0576:phyxminiproblems-problemphyxmini0576-timeinseconds}
  \lean{PhyXMiniProblems.ProblemPhyXMini0576.timeInSeconds}
  \uses{def:physics:phyx-mini-0576:phyxminiproblems-problemphyxmini0576-timequantity, def:physics:phyx-mini-0576:phyxminiproblems-problemphyxmini0576-timereadout}
  Seconds readout of a physical elapsed time
\end{definition}
\begin{proof}
  This is the declared model definition of time In Seconds.
\end{proof}
\begin{definition}[decay Constant In Per Second]
  \label{def:physics:phyx-mini-0576:phyxminiproblems-problemphyxmini0576-decayconstantinpersecond}
  \lean{PhyXMiniProblems.ProblemPhyXMini0576.decayConstantInPerSecond}
  \uses{def:physics:phyx-mini-0576:phyxminiproblems-problemphyxmini0576-decayconstantquantity, def:physics:phyx-mini-0576:phyxminiproblems-problemphyxmini0576-decayconstantreadout}
  Per-second readout of a radioactive decay constant
\end{definition}
\begin{proof}
  This is the declared model definition of decay Constant In Per Second.
\end{proof}
\begin{definition}[activity In Becquerels]
  \label{def:physics:phyx-mini-0576:phyxminiproblems-problemphyxmini0576-activityinbecquerels}
  \lean{PhyXMiniProblems.ProblemPhyXMini0576.activityInBecquerels}
  \uses{def:physics:phyx-mini-0576:phyxminiproblems-problemphyxmini0576-activityquantity, def:physics:phyx-mini-0576:phyxminiproblems-problemphyxmini0576-activityreadout}
  Becquerel readout of an activity; one becquerel is one decay per second
\end{definition}
\begin{proof}
  This is the declared model definition of activity In Becquerels.
\end{proof}
\begin{definition}[Radioactive Sample]
  \label{def:physics:phyx-mini-0576:phyxminiproblems-problemphyxmini0576-radioactivesample}
  \lean{PhyXMiniProblems.ProblemPhyXMini0576.RadioactiveSample}
  \uses{def:physics:phyx-mini-0576:phyxminiproblems-problemphyxmini0576-timequantity, def:physics:phyx-mini-0576:phyxminiproblems-problemphyxmini0576-decayconstantquantity, def:physics:phyx-mini-0576:phyxminiproblems-problemphyxmini0576-activityquantity}
  An idealized radioactive sample. Its expected parent population and activity at a time are independent observables here; the governing laws below relate them without defining either observable to be the requested answer
\end{definition}
\begin{proof}
  This is the declared model definition of Radioactive Sample.
\end{proof}
\begin{definition}[Decay Curve Color]
  \label{def:physics:phyx-mini-0576:phyxminiproblems-problemphyxmini0576-decaycurvecolor}
  \lean{PhyXMiniProblems.ProblemPhyXMini0576.DecayCurveColor}
  The color of the curve visible in the supplied bitmap
\end{definition}
\begin{proof}
  This is the declared model definition of Decay Curve Color.
\end{proof}
\begin{definition}[Decay Curve Trend]
  \label{def:physics:phyx-mini-0576:phyxminiproblems-problemphyxmini0576-decaycurvetrend}
  \lean{PhyXMiniProblems.ProblemPhyXMini0576.DecayCurveTrend}
  Qualitative vertical trend of the plotted parent population
\end{definition}
\begin{proof}
  This is the declared model definition of Decay Curve Trend.
\end{proof}
\begin{definition}[Decay Curve Shape]
  \label{def:physics:phyx-mini-0576:phyxminiproblems-problemphyxmini0576-decaycurveshape}
  \lean{PhyXMiniProblems.ProblemPhyXMini0576.DecayCurveShape}
  Physical model represented by the shape of the plotted curve
\end{definition}
\begin{proof}
  This is the declared model definition of Decay Curve Shape.
\end{proof}
\begin{definition}[Parent Decay Figure]
  \label{def:physics:phyx-mini-0576:phyxminiproblems-problemphyxmini0576-parentdecayfigure}
  \lean{PhyXMiniProblems.ProblemPhyXMini0576.ParentDecayFigure}
  \uses{def:physics:phyx-mini-0576:phyxminiproblems-problemphyxmini0576-decaycurvecolor, def:physics:phyx-mini-0576:phyxminiproblems-problemphyxmini0576-decaycurvetrend, def:physics:phyx-mini-0576:phyxminiproblems-problemphyxmini0576-decaycurveshape}
  Literal labels and grid information read from the primary image. Grid indices are normalized diagram coordinates, not physical times or parent counts
\end{definition}
\begin{proof}
  This is the declared model definition of Parent Decay Figure.
\end{proof}
\begin{definition}[Radioactive Decay Setup]
  \label{def:physics:phyx-mini-0576:phyxminiproblems-problemphyxmini0576-radioactivedecaysetup}
  \lean{PhyXMiniProblems.ProblemPhyXMini0576.RadioactiveDecaySetup}
  \uses{def:physics:phyx-mini-0576:phyxminiproblems-problemphyxmini0576-timequantity, def:physics:phyx-mini-0576:phyxminiproblems-problemphyxmini0576-radioactivesample, def:physics:phyx-mini-0576:phyxminiproblems-problemphyxmini0576-parentdecayfigure}
  The sample, the two axis scales, the queried time, and the half-life point read from the graph. Each physical quantity remains independent of its scalar readout and of the answer choices
\end{definition}
\begin{proof}
  This is the declared model definition of Radioactive Decay Setup.
\end{proof}
\begin{definition}[Matches Problem Readouts]
  \label{def:physics:phyx-mini-0576:phyxminiproblems-problemphyxmini0576-matchesproblemreadouts}
  \lean{PhyXMiniProblems.ProblemPhyXMini0576.MatchesProblemReadouts}
  \uses{def:physics:phyx-mini-0576:phyxminiproblems-problemphyxmini0576-timeinseconds, def:physics:phyx-mini-0576:phyxminiproblems-problemphyxmini0576-radioactivedecaysetup}
  Numerical scales and query time stated alongside the graph
\end{definition}
\begin{proof}
  This is the declared model definition of Matches Problem Readouts.
\end{proof}
\begin{definition}[Matches Supplied Parent Decay Figure]
  \label{def:physics:phyx-mini-0576:phyxminiproblems-problemphyxmini0576-matchessuppliedparentdecayfigure}
  \lean{PhyXMiniProblems.ProblemPhyXMini0576.MatchesSuppliedParentDecayFigure}
  \uses{def:physics:phyx-mini-0576:phyxminiproblems-problemphyxmini0576-timeinseconds, def:physics:phyx-mini-0576:phyxminiproblems-problemphyxmini0576-radioactivedecaysetup}
  Primary-image evidence: the axes are labelled t and N, the endpoint scales are t\_s and N\_s, and a magenta decreasing exponential occupies a ten-by-ten grid. The curve crosses half of N\_s at horizontal grid column two and vertical row five. The last two fields interpret that calibrated crossing as the sample's half-population point; neither field concerns the requested activity at 27 s
\end{definition}
\begin{proof}
  This is the declared model definition of Matches Supplied Parent Decay Figure.
\end{proof}
\begin{definition}[Has Physical Radioactive Parameters]
  \label{def:physics:phyx-mini-0576:phyxminiproblems-problemphyxmini0576-hasphysicalradioactiveparameters}
  \lean{PhyXMiniProblems.ProblemPhyXMini0576.HasPhysicalRadioactiveParameters}
  \uses{def:physics:phyx-mini-0576:phyxminiproblems-problemphyxmini0576-timeinseconds, def:physics:phyx-mini-0576:phyxminiproblems-problemphyxmini0576-decayconstantinpersecond, def:physics:phyx-mini-0576:phyxminiproblems-problemphyxmini0576-radioactivedecaysetup}
  Positive, nondegenerate parameters of the radioactive-decay scenario
\end{definition}
\begin{proof}
  This is the declared model definition of Has Physical Radioactive Parameters.
\end{proof}
\begin{definition}[Satisfies Radioactive Decay Laws]
  \label{def:physics:phyx-mini-0576:phyxminiproblems-problemphyxmini0576-satisfiesradioactivedecaylaws}
  \lean{PhyXMiniProblems.ProblemPhyXMini0576.SatisfiesRadioactiveDecayLaws}
  \uses{def:physics:phyx-mini-0576:phyxminiproblems-problemphyxmini0576-timequantity, def:physics:phyx-mini-0576:phyxminiproblems-problemphyxmini0576-timeinseconds, def:physics:phyx-mini-0576:phyxminiproblems-problemphyxmini0576-decayconstantinpersecond, def:physics:phyx-mini-0576:phyxminiproblems-problemphyxmini0576-activityinbecquerels, def:physics:phyx-mini-0576:phyxminiproblems-problemphyxmini0576-radioactivedecaysetup}
  The standard exponential parent-population law and A(t) = λ N(t). Both are general governing laws quantified over every elapsed time. In particular, they do not state a numerical activity at the queried time
\end{definition}
\begin{proof}
  This is the declared model definition of Satisfies Radioactive Decay Laws.
\end{proof}
\begin{definition}[Answer Choice]
  \label{def:physics:phyx-mini-0576:phyxminiproblems-problemphyxmini0576-answerchoice}
  \lean{PhyXMiniProblems.ProblemPhyXMini0576.AnswerChoice}
  The four answer labels printed in the source
\end{definition}
\begin{proof}
  This is the declared model definition of Answer Choice.
\end{proof}
\begin{definition}[displayed Activity In Becquerels]
  \label{def:physics:phyx-mini-0576:phyxminiproblems-problemphyxmini0576-displayedactivityinbecquerels}
  \lean{PhyXMiniProblems.ProblemPhyXMini0576.displayedActivityInBecquerels}
  \uses{def:physics:phyx-mini-0576:phyxminiproblems-problemphyxmini0576-answerchoice}
  Becquerel value printed beside each answer label
\end{definition}
\begin{proof}
  This is the declared model definition of displayed Activity In Becquerels.
\end{proof}
\begin{definition}[whole Becquerel Rounding Tolerance]
  \label{def:physics:phyx-mini-0576:phyxminiproblems-problemphyxmini0576-wholebecquerelroundingtolerance}
  \lean{PhyXMiniProblems.ProblemPhyXMini0576.wholeBecquerelRoundingTolerance}
  The answer choices are reported to the nearest whole becquerel
\end{definition}
\begin{proof}
  This is the declared model definition of whole Becquerel Rounding Tolerance.
\end{proof}
\begin{definition}[Is Correct Rounded Answer Choice]
  \label{def:physics:phyx-mini-0576:phyxminiproblems-problemphyxmini0576-iscorrectroundedanswerchoice}
  \lean{PhyXMiniProblems.ProblemPhyXMini0576.IsCorrectRoundedAnswerChoice}
  \uses{def:physics:phyx-mini-0576:phyxminiproblems-problemphyxmini0576-activityinbecquerels, def:physics:phyx-mini-0576:phyxminiproblems-problemphyxmini0576-radioactivedecaysetup, def:physics:phyx-mini-0576:phyxminiproblems-problemphyxmini0576-answerchoice, def:physics:phyx-mini-0576:phyxminiproblems-problemphyxmini0576-displayedactivityinbecquerels, def:physics:phyx-mini-0576:phyxminiproblems-problemphyxmini0576-wholebecquerelroundingtolerance}
  A displayed choice agrees with the physical prediction when the exact model activity lies within half a becquerel of the displayed whole-number value
\end{definition}
\begin{proof}
  This is the declared model definition of Is Correct Rounded Answer Choice.
\end{proof}
% --- Archon declaration topology end ---
% --- Archon physics formalization source end ---
